% ch10_d 第10章 §10.5–§10.6 (书页 479–491 / p494–p506)
% 实际覆盖范围（以 PNG 为准）：本块不含 §10.5——§10.5、§10.6.1 及 §10.6.2 前半
% （书页 466–478）属 ch10_c 块；本块自 §10.6.2 中段（书页 479 顶）起，
% 至 §10.7.3 中段（书页 491 末）止。
% 块界说明：
%   起点——书页 479（p494）顶部 "In the following, we are going to show that..."
%   为新起段落，书页 478（p493）末句 "...the two nearest-neighbor rotors around
%   the i site." 已完整结束，故本块不设 %JOIN%，亦无 ch10_c 收尾内容可写入
%   %--C-TAIL-- 区。
%   终点——书页 491（p506）末句 "In the continuum limit, ... and B is the
%   magnetic field" 延续到书页 492（p507，ch10_e 块）顶部 "strength of the
%   U(1) gauge field. Comparing this with..."。按约定，半句译文（"……而 $\pmb B$ 是"）
%   留在本文件正文最末，未写入任何 TAIL 区；请 ch10_e 以 %JOIN% 起头续译
%   "该 $U(1)$ 规范场的磁场强度。把它与……"，勿重复勿遗漏。

下面我们将证明，在 $U\gg t$、$J$ 且 $t^{2}/U\gg J$ 的极限下，上述二维和三维模型都包含一个低能集体模式。由于其非局域的性质，这种集体模式与自旋波、声子这类通常的集体模式非常不同。实际上，它不对应于任何局域序参量的涨落，必须用 $U(1)$ 规范理论来描写。我们把这样的集体模式称为 $U(1)$ 规范涨落。该集体模式在二维模型中具有指数式小的能隙 $\Delta\sim\mathrm{e}^{-C\sqrt{t^{2}/JU}}$，在三维模型中则严格无能隙。三维模型中的集体模式在各个方面都表现得像我们宇宙中的光。因此，我们也可以称它为人工光。

上述模型引人注目之处在于，它们在格子尺度上没有任何规范结构的迹象。$U(1)$ 规范涨落及与之相联系的规范结构是在低能下演生的。这提示了一种可能：我们宇宙中的光或许也以类似的方式出现。规范结构可能会在强关联量子系统中自然地演生。本节要研究的转子模型是包含人工光的最简单模型。下面为简单起见，我们将集中于二维转子模型。计算与结果可以容易地推广到三维转子模型。

为了理解我们二维转子系统的动力学，让我们先设 $J=t=0$ 且 $U>0$。此时，哈密顿量由执行局域投影的对易项构成。基态高度简并，构成一个低能子空间。其中一个基态是每个转子的 $S^{z}_{\pmb i}=0$ 的态。其他基态可以从这个基态出发来构造：在正方格子上画一条有向回路，然后沿回路交替地把转子的角动量增加或减少 1（见图 10.12）。重复这一过程可以构造出所有简并基态。我们看到，投影空间中的涨落由回路表示。尽管低能子空间是通过局域投影得到的，它却带有一些非局域的特征。若 $t$ 与 $J$ 非零，则 $t$ 项会使这些回路涨落，而 $J$ 项会给这些回路正比于回路长度的能量。显然，当 $U\gg J,t$ 时，系统的低能性质由回路的涨落决定。系统可以有两个不同的相。当 $J\gg t$ 时，系统只有少许小回路涨落。当 $J\ll t$ 时，系统有许多大回路涨落，它们甚至可以充满整个空间。后面我们会看到，这些大回路涨落实际上正是 $U(1)$ 规范涨落。我们注意到，回路可以相交与交叠。实际上，一个典型的回路看起来更像一张闭合弦的网。下面，为强调这种涨落的分支结构，我们把回路涨落称为闭合弦网涨落。

简并基态在由
\begin{equation*}
U(\phi_{\pmb i})=\mathrm{e}^{\mathrm{i}\sum_{\pmb i}((-)^{\pmb i}\phi_{\pmb i}\sum_{\alpha}S^{z}_{\pmb i+\alpha})}\tag{10.6.8}
\end{equation*}
生成的局域对称变换下不变，其中 $(-)^{\pmb i}=(-)^{i_x+i_y}$。上述变换正是规范变换。因此，我们也可以说简并基态是规范不变的。

用与 10.6.1 小节相似的计算，我们发现我们的格子转子模型在大 $U$ 极限下可以用如下低能有效拉格朗日量描写：
\begin{equation*}
L=\frac{1}{4J}\sum_{\pmb i,\pmb\mu=\pmb x,\pmb y}[\dot{a}_{\pmb i,\pmb i+\pmb\mu}+a_0(\pmb i)-a_0(\pmb i+\pmb\mu)]^{2}+g\sum_{\pmb i}\cos(\Phi_{\pmb i})\tag{10.6.9}
\end{equation*}
这里 $a_{\pmb i,\pmb i+\pmb x}=(-)^{\pmb i}\theta_{\pmb i+\frac{\pmb x}{2}}$，$a_{\pmb i,\pmb i+\pmb y}=(-)^{\pmb i}\theta_{\pmb i+\frac{\pmb y}{2}}$，而 $a_{\pmb i\pmb j}=-a_{\pmb j\pmb i}$。另外，$\Phi_{\pmb i}$ 是通过 $\pmb i$ 处方格的 $U(1)$ 通量，即 $\Phi_{\pmb i}=a_{\pmb i,\pmb i+\pmb x}+a_{\pmb i+\pmb x,\pmb i+\pmb x+\pmb y}+a_{\pmb i+\pmb x+\pmb y,\pmb i+\pmb y}+a_{\pmb i+\pmb y,\pmb i}$。推广 10.6.1 小节的二阶微扰计算，我们得到
\begin{equation*}
g=2(t_1t_3+t_2t_4)/U.
\end{equation*}

式 (10.6.9) 描写一个 $U(1)$ 格点规范理论。

这里我们想强调，具有 $U(1)$ 规范理论形式的拉格朗日量 (10.6.9) 并不一定意味着低能激发的行为像规范玻色子。这是因为量子涨落可以非常强，来自拉格朗日量的直观图像可能完全造成误导。要有行为像规范玻色子的激发，我们需要选取 $J$ 与 $g$，使拉格朗日量 (10.6.9) 处于半经典极限。为了看清何时能够达到半经典极限，我们把格点规范理论的作用量写成如下无量纲形式：
\begin{equation*}
S=\sqrt{\frac{g}{J}}\int\mathrm{d}\tilde t\left(\frac{1}{4}\sum_{\pmb i,\pmb\mu=\pmb x,\pmb y}[\partial_{\tilde t}a_{\pmb i,\pmb i+\pmb\mu}+\tilde a_0(\pmb i)-\tilde a_0(\pmb i+\pmb\mu)]^{2}+\sum_{\pmb i}\cos(\Phi_{\pmb i})\right)
\end{equation*}
其中 $\tilde t=\sqrt{gJ}\,t$，$\tilde a_0=a_0/\sqrt{gJ}$。我们发现，当 $g/J=t^{2}/JU\gg1$ 时即达到半经典极限。在这一极限中，我们的模型包含 $U(1)$ 规范玻色子（即人工光）作为其唯一的低能集体激发。

由于瞬子效应，$U(1)$ 规范激发在 $1+2$ 维中会产生能隙（Polyakov, 1977）（见 6.3.2 小节）。瞬子效应与一个方格上 $U(1)$ 通量 $\Phi$ 从 0 到 $2\pi$ 的改变相联系。为了估计瞬子效应的重要性，让我们考虑只含单个方格的模型（即前面讨论过的四转子模型）。这样的模型由式 (6.4.11) 描写。瞬子效应对应于一条路径 $\Phi(t)$（依赖时间的通量），其中 $\Phi$ 从 $\Phi(-\infty)=0$ 变到 $\Phi(+\infty)=2\pi$。为了估计瞬子作用量，我们设
\begin{equation*}
\Phi(t)=\left\{\begin{array}{ll}0,&\text{当 } t<0\\ 2\pi t/T,&\text{当 } 0<t<T\\ 2\pi,&\text{当 } T<t\end{array}\right.
\end{equation*}
瞬子作用量的最小值为
\begin{equation*}
S_c=\pi\sqrt{g/J}
\end{equation*}
此时 $T=\pi/2\sqrt{gJ}$。由瞬子气体在时间轴上的密度 $\sqrt{Jg}\,\mathrm{e}^{-S_c}$，我们估计 $U(1)$ 规范玻色子的能隙为
\begin{equation*}
\Delta_{\mathrm{gauge}}\sim\sqrt{Jg}\,\mathrm{e}^{-\pi\sqrt{g/J}}\tag{10.6.10}
\end{equation*}

我们看到，当 $g/J\gg1$ 时，能隙可以非常小，低能涨落非常像无能隙的光子。当然，真正无能隙的人工光只存在于 $(3+1)$ 维模型中，例如三维转子模型 (10.6.7)。

三维转子模型在大 $U$ 极限下具有如下低能有效拉格朗日量：
\begin{equation*}
L=\frac{1}{4J}\sum_{\pmb i,\pmb\mu=\pmb x,\pmb y,z}[\dot{a}_{\pmb i,\pmb i+\pmb\mu}+a_0(\pmb i)-a_0(\pmb i+\pmb\mu)]^{2}+\sum_{\pmb p}g_{\pmb p}\cos(\Phi_{\pmb p})\tag{10.6.11}
\end{equation*}
这里 $\sum_{\pmb p}$ 对所有方格求和，$\Phi_{\pmb p}$ 是通过方格 $\pmb p$ 的 $U(1)$ 通量。另外，$g_{\pmb p}$ 可以依赖方格的取向，其数值可通过式 (10.6.7) 中的 $t_{\langle\pmb\alpha_1\pmb\alpha_2\rangle}$ 调节。

\subsection{人工光与人工电荷的弦网理论}

\begin{keypoints}
\item 规范理论是披着伪装的闭合弦网理论。
\item 禁闭相是具有稀疏小弦的相；库仑相是具有充满整个空间的大闭合弦网的相。
\item 规范玻色子（在库仑相中）对应于大闭合弦网的涨落，而规范电荷是开弦的端点。
\end{keypoints}

如前所述，我们的转子模型中 $U$ 以下的低能激发由 $S^{z}$ 增/减的闭合弦网描写。为了把这个图像弄得更精确，我们想在格子上定义一个弦网理论。

弦网理论的希尔伯特空间是我们转子模型 (10.6.6) 的希尔伯特空间的一个子空间。为了构造弦网希尔伯特空间，我们首先需要引入一个弦产生算符。弦产生算符由 $\mathrm{e}^{\pm\mathrm{i}\theta_{\langle\pmb i\pmb j\rangle}}$ 算符的乘积构成：
%FIG{10.13}: {(a) 闭合弦网；(b) 开弦网。}
\begin{equation*}
U(C)=\prod_{\langle\pmb i\pmb j\rangle}\mathrm{e}^{\mathrm{i}(-)^{\pmb i}\theta_{\langle\pmb i\pmb j\rangle}}\tag{10.6.12}
\end{equation*}
其中 $C$ 是连接正方格子上最近邻格点的一根弦，乘积 $\prod_{\langle\pmb i\pmb j\rangle}$ 遍历构成该弦的所有最近邻链 $\langle\pmb i\pmb j\rangle$，而 $\theta_{\langle\pmb i\pmb j\rangle}$ 是链 $\langle\pmb i\pmb j\rangle$ 上转子的 $\theta$ 场。

由于弦可以相交与交叠，弦看起来更像一张网。我们把 $C$ 称为弦网，把 $U(C)$ 称为弦网算符。若 $C$ 由一组闭合弦构成，我们就称它为闭合弦网；若 $C$ 至少含有一根开弦，就称它为开弦网（见图 10.13）。

弦网希尔伯特空间包含一个所有 $S^{z}_{\pmb I}=0$ 的态。按定义，这样的态对应于无弦态。如果我们把闭合弦算符 (10.6.12) 作用到 $S^{z}_{\pmb I}=0$ 态上，就得到弦网希尔伯特空间中的另一个态。这样的态由沿闭合弦 $C$ 的 $S^{z}_{\pmb I}=\pm1$ 构成，可以视为含一根闭合弦 $C$ 的态。弦网希尔伯特空间中的其他态对应于多弦态，可以通过反复把闭合弦算符 (10.6.12) 作用到 $S^{z}_{\pmb I}=0$ 态上产生。

我们弦网理论的哈密顿量由限制在弦网子空间中的转子哈密顿量 (10.6.6) 得到，其形式为
\begin{equation*}
H_{str}=\sum_{\langle\pmb i\pmb j\rangle}J(S^{z}_{\pmb i\pmb j})^{2}-\sum_{\pmb p}\frac{1}{2}(gW_{\pmb p}+h.c.)\tag{10.6.13}
\end{equation*}
其中求和 $\sum_{\langle\pmb i\pmb j\rangle}$ 遍历正方格子上所有最近邻链 $\langle\pmb i\pmb j\rangle$，$S^{z}_{\pmb i\pmb j}$ 是最近邻链 $\langle\pmb i\pmb j\rangle$ 上转子的角动量算符，$\sum_{\pmb p}$ 遍历正方格子的所有方格，而 $W_{\pmb p}$ 是绕方格 $\pmb p$ 的闭合弦的闭合弦算符。可以验证，上述哈密顿量确实作用在弦网希尔伯特空间之内。$J$ 项赋予弦网中的弦有限的弦张力，而 $g$ 项作为弦的“跳跃”项使弦网发生涨落。式 (10.6.13) 中的 $J$ 项直接来自式 (10.6.6) 中的 $J$ 项，而 $g$ 项来自式 (10.6.6) 中的 $t$ 项的二阶微扰展开。

由构造可以清楚地看到，弦网希尔伯特空间与我们的模型 (10.6.6) 的低能希尔伯特空间（由能量小于 $U$ 的态构成）恒同。由有效格点规范理论 (10.6.9) 的推导也同样清楚，弦网哈密顿量 (10.6.13) 与格点规范拉格朗日量 (10.6.9) 直接相关。实际上，格点规范理论的哈密顿量与弦网哈密顿量 (10.6.13) 恒同。弦网理论中的 $\sum_{\pmb i\pmb j}J(S^{z}_{\pmb i\pmb j})^{2}$ 项对应于规范理论中的 $\frac{1}{4J}\sum_{\langle\pmb i\pmb j\rangle}[\dot{a}_{\pmb i\pmb j}+a_0(\pmb i)-a_0(\pmb j)]^{2}$ 项，弦网理论中的 $\sum_{\pmb p}\frac{1}{2}g(W_{\pmb p}+h.c.)$ 项对应于规范理论中的 $g\sum_{\pmb p}\cos(\Phi_{\pmb p})$ 项。由于 $S^{z}\sim J^{-1}\dot{\theta}_{\langle\pmb i\pmb j\rangle}=J^{-1}\dot{a}_{\pmb i\pmb j}$ 对应于沿链的电通量，$S^{z}$ 增/减的闭合弦便对应于电通量管的一个回路。我们想强调，规范理论与弦网理论之间的上述联系不仅适用于二维模型 (10.6.6)，也适用于任何维数中的模型。

我们看到，$U(1)$ 规范理论 (10.6.9) 实际上是闭合弦网的一个动力学理论。通常人们期望闭合弦网的动力学理论像 (10.6.13) 那样用弦网写成。然而，由于我们更熟悉场论，前几节所做的事情可以视为用场论描写弦网理论的一种尝试。借助某些数学技巧，我们达到了目的：能够把弦网理论写成规范场论的形式。规范场论是一类特殊的场论，其中的场\emph{并不}对应于物理自由度，而且其物理希尔伯特空间是非局域的（即总物理希尔伯特空间不能写成局域希尔伯特空间的直积）。我们在这里想说明的要点是：规范理论（至少这里讨论的规范理论）是一种披着伪装的闭合弦网理论。换言之，规范理论与闭合弦网理论互为对偶。实际上，如果把时间离散化、考虑一个时空格子，我们就能找到 $U(1)$ 格点规范理论与时空中的膜的统计模型之间的精确映射（Banks et al., 1977; Savit, 1980）。

在大 $J/g$ 极限（因而是大 $\Delta_{\mathrm{gauge}}$ 极限）下，转子模型与弦网模型的基态都是每个转子 $S^{z}=0$ 的态。在这个相中，闭合弦网（或电通量管）涨落不大，能量正比于其长度。这意味着 $U(1)$ 规范理论处于禁闭相。在小 $J/g$ 极限下，闭合弦网强烈涨落，空间充满任意尺寸的闭合弦网。根据上一小节的计算，我们注意到小 $J/g$ 相也可以视为具有无能隙规范玻色子的库仑相。把这两幅图像结合起来，我们看到无能隙规范玻色子对应于大闭合弦网的涨落。

在把闭合弦网与人工光联系起来之后，我们现在转向人工电荷（artificial charge）。要为人造 $U(1)$ 规范场产生一对带相反人工电荷的粒子，我们需要画一根开弦，并沿弦交替地增减转子的 $S^{z}$（见图 10.12）。开弦的端点，作为电通量管的端点，对应于带相反人工电荷的粒子。我们注意到，与转子不同，带电粒子生活在正方格子的格点上。在禁闭相中，连接两个人工电荷的弦涨落不大，弦的能量正比于弦的长度。我们看到，人工电荷之间存在线性禁闭。

在小 $J/g$ 极限下，大的 $g$ 引起闭合弦网的强烈涨落，从而导致无能隙的 $U(1)$ 规范涨落。连接两个电荷的弦的强烈涨落还把线性禁闭势变为电荷之间的 $\log(r)$ 势。

为了理解带人工电荷的粒子的动力学并导出 $\log(r)$ 势，让我们导出这些带电粒子的低能有效理论。让我们先设 $J=t=0$。把开弦算符 (10.6.12) 作用到基态上，可以产生一对带相反单位人工电荷的带电粒子。在开弦的端点处，$(\sum_{\alpha}S^{z}_{\pmb i+\alpha})^{2}=1$。我们发现每个带电粒子具有能量 $U$，它来自 $U\sum_{\pmb i}(\sum_{\alpha}S^{z}_{\pmb i+\alpha})^{2}$ 这一项。弦本身不耗费能量。如果把同一开弦算符作用 $n$ 次，就在开弦的端点产生 $n$ 个单位的相反人工电荷。$n$ 个单位电荷的能量是 $n^{2}U$。

含电荷的规范理论包含两部分，即电荷与弦。让我们先只考虑电荷，把带电粒子当作独立粒子处理。此时，带电粒子的总希尔伯特空间由态 $|\{n_{\pmb i}\}\rangle$ 构成，其中 $n_{\pmb i}$ 是正方格子格点 $\pmb i$ 上人工电荷的数目。这里 $|\{n_{\pmb i}\}\rangle$ 是一个能量本征态，能量为 $E=U\sum_{\pmb i}n^{2}_{\pmb i}$。这样的系统可以用拉格朗日量
\begin{equation*}
L=\sum_{\pmb i}\frac{1}{4U}\dot{\varphi}^{2}_{\pmb i}\tag{10.6.14}
\end{equation*}
描写，其中 $\varphi_{\pmb i}$ 是一个角变量。产生单位电荷带电粒子的算符是 $\mathrm{e}^{\mathrm{i}\varphi_{\pmb i}}$。

规范涨落由弦和式 (10.6.9) 描写。在 $U(1)$ 规范理论中，基态是规范不变的（在远离电荷处）。由式 (10.6.8) 可见，规范不变性意味着 $\sum_{\alpha}S^{z}_{\pmb i+\alpha}=0$。因此在远离电荷处没有弦的开端。

为了找到电荷与规范场的耦合理论，我们把带电粒子总是处在开弦端点这一事实包括进来。换言之，规范理论的物理态包含开弦，但开弦只能终止在电荷上。这一约束可以由规范不变性施加。于是我们可以利用规范不变性，把电荷拉格朗日量 (10.6.14) 与规范拉格朗日量 (10.6.9) 合在一起，写出耦合理论。利用规范不变性，我们发现组合拉格朗日量具有如下形式（见问题 6.4.5）：
\begin{equation*}
L=\sum_{\pmb i}\frac{(\dot{\varphi}_{\pmb i}+a_0(\pmb i))^{2}}{4U}+\sum_{\pmb i,\pmb\mu=\pmb x,\pmb y}\frac{[\dot{a}_{\pmb i,\pmb i+\pmb\mu}+a_0(\pmb i)-a_0(\pmb i+\pmb\mu)]^{2}}{4J}+g\sum_{\pmb i}\cos(\Phi_{\pmb i})
\end{equation*}
把规范场包括进来之后，单电荷产生算符 $\mathrm{e}^{\mathrm{i}\varphi_{\pmb i}}$ 不再是物理的，因为它不是规范不变的。规范不变的算符
\begin{equation*}
\mathrm{e}^{-\mathrm{i}\varphi_{\pmb i_1}}\mathrm{e}^{\mathrm{i}a_{\pmb i_1\pmb i_2}}\mathrm{e}^{\mathrm{i}a_{\pmb i_2\pmb i_3}}\dots\mathrm{e}^{\mathrm{i}a_{\pmb i_{N-1}\pmb i_N}}\mathrm{e}^{\mathrm{i}\varphi_{\pmb i_N}}
\end{equation*}
总是产生一对相反电荷，连同连接两电荷的开弦。因此，开弦总是终止在电荷上。实际上，上述规范不变算符正是开弦算符 (10.6.12)。我们还看到，弦算符 (10.6.12) 与 Wegner--Wilson 圈算符密切相关（Wegner, 1971; Wilson, 1974; Kogut, 1979）。

$t$ 项产生带电粒子在正方格子上次近邻之间的跳跃。因此，若 $t\neq0$，带电粒子将具有非平庸的色散。相应的拉格朗日量为
\begin{equation*}
L=\sum_{\pmb i}\frac{(\dot{\varphi}_{\pmb i}+a_0(\pmb i))^{2}}{4U}-\sum_{\langle\pmb i\pmb j\rangle,a=1,2}t(\mathrm{e}^{\mathrm{i}(\varphi_{\pmb i}-\varphi_{\pmb j}-a_{\pmb i\pmb k_a}-a_{\pmb k_a\pmb j})}+h.c.)\tag{10.6.15}
\end{equation*}
其中 $\langle\pmb i\pmb j\rangle$ 是正方格子上的次近邻，$\pmb k_{1,2}$ 是格点 $\pmb i$ 与 $\pmb j$ 之间的两个格点。上述拉格朗日量还告诉我们，带电粒子是玻色子。

我们还注意到，$S^{z}_{\pmb i\pmb j}$ 的增加对应于 $\pmb i$ 与 $\pmb j$ 处的两个人工电荷。因此，每一单位人工电荷携带 1/2 角动量！（注意，总角动量 $\sum_{\langle\pmb i\pmb j\rangle}S^{z}_{\langle\pmb i\pmb j\rangle}$ 是守恒量。）

\subsection{二维与三维转子系统的物理性质}

\begin{keypoints}
\item 格点规范理论的连续极限。
\item 人工光的速度与精细结构常数。
\end{keypoints}

为了理解二维模型中人工光的物理性质，让我们按如下方式取式 (10.6.9) 的连续极限：
\begin{equation*}
a_{\pmb i\pmb j}=\delta\pmb x_{\langle\pmb i\pmb j\rangle}\cdot\pmb a(\pmb x),\qquad a_0(\pmb i)=a_0(\pmb x),
\end{equation*}
其中 $\pmb a=(a_x,a_y)$ 是一个二维矢量场（二维中的矢量规范势），$a_0$ 对应于势场，$\pmb x$ 取在格点 $\pmb i$ 附近，$\delta\pmb x_{\langle\pmb i\pmb j\rangle}$ 是连接正方格子上 $\pmb i$ 与 $\pmb j$ 两格点的矢量，而 $l$ 是正方格子的晶格常数。在连续极限下，拉格朗日量 (10.6.9) 变为
\begin{equation*}
L=\int\mathrm{d}^{2}\pmb x\ \left(\frac{1}{4J}\pmb e^{2}-\frac{gl^{2}}{2}\pmb b^{2}\right)\tag{10.6.16}
\end{equation*}
其中 $\pmb e=\partial_t\pmb a-\partial_{\pmb x}a_0$ 与 $\pmb b=\partial_x a_y-\partial_y a_x$ 分别是相应的仿造电场与仿造磁场。我们看到，人工光的速度为 $c_a=\sqrt{2gJl^{2}/h^{2}}$，人工光的带宽约为 $E_a=2c_a\hbar/l=\sqrt{8gJ}$。

由式 (10.6.15)，我们得到描写二维模型中带电粒子（在 $U\gg t$ 极限下）的连续拉格朗日量：
\begin{align*}
L=\int\mathrm{d}^{2}\pmb x\ \sum_{I=1,2}\Big(\phi^{\dagger}_{I}&(\mathrm{i}\partial_t-a_0-U)\phi_I-\frac{8tl^{2}}{2}\big|(\partial_{i}+\mathrm{i}a_{i})\phi_I\big|^{2}\\
+\bar{\phi}^{\dagger}_{I}&(\mathrm{i}\partial_t+a_0-U)\bar{\phi}_I-\frac{8tl^{2}}{2}\big|(\partial_{i}-\mathrm{i}a_{i})\bar{\phi}_I\big|^{2}\Big)
\end{align*}
（译注：原文两处时间导数下标印作 $\partial_l$，应为 $\partial_t$，此处已按本意译出。）

其中 $\phi_I$ 描写带正电的玻色子，$\bar{\phi}_I$ 描写带负电的玻色子，$\psi_1$ 与 $\bar{\psi}_1$ 描写正方格子偶格点上带电的玻色子，$\psi_2$ 与 $\bar{\psi}_2$ 描写奇格点上带电的玻色子。产生一对带电玻色子耗费能量 $2U$。玻色子的质量为 $m=(8tl^{2})^{-1}$，且 $mc_a^{2}=gJ/4t$。我们想指出，玻色子速度 $\sim 8tl$ 可以大于人工光的速度。正电荷与负电荷之间的势能为 $V(r)=\frac{2J}{\pi}\ln r$。

对于三维模型，连续极限下的拉格朗日量为
\begin{equation*}
L=\int\mathrm{d}^{3}\pmb x\ \left(\frac{1}{4Jl}\pmb e^{2}-\frac{gl}{2}\pmb b^{2}\right)\tag{10.6.17}
\end{equation*}
其中 $\pmb e$ 与 $\pmb b$ 分别是三维中的仿造电场与仿造磁场。人工光的速度为 $c_a=\sqrt{2gJl^{2}/h^{2}}$。上述三维拉格朗日量可以改写成更标准的形式
\begin{equation*}
L=\int\mathrm{d}^{3}\pmb x\ \frac{1}{8\pi\alpha}\left(\frac{1}{c_a}\pmb e^{2}-c_a\pmb b^{2}\right)\tag{10.6.18}
\end{equation*}
其中 $\alpha=\frac{1}{2\pi}\sqrt{J/2g}$ 是人工精细结构常数。带电玻色子的质量为 $m=(8tl^{2})^{-1}$。人工原子（artificial atom，即一个带正电玻色子与一个带负电玻色子的束缚态）具有 $\frac{1}{4}mc_a^{2}\alpha^{2}$ 量级的能级间距和 $2/\alpha mc_a$ 量级的尺寸。

\subsection*{问题 10.6.2}

验证式 (10.6.16) 与 (10.6.17)。

\section{从 $SU(N_f)$ 自旋模型演生的光与电子}

\begin{keypoints}
\item 局域玻色模型中非平庸的量子序（即弦网凝聚）可以同时导致无质量的光子与无质量的费米子。弦网凝聚提供了统一光与费米子的一条途径。
\item 光子与费米子的无质量性由刻画弦网凝聚的 PSG 保护。
\end{keypoints}

正如本章开头所强调的，理解光与电子的存在，就是理解这些粒子为什么是无质量的（与 Planck 质量相比近乎无质量）。对于一个一般的相互作用系统，无质量（即无能隙）激发是罕见的。如果它们存在，那么它们的存在必有原因。

我们知道，对称性破缺能够产生并保护无能隙的 Nambu--Goldstone 模。9.10 节中提出，除对称性破缺之外，量子序及与之相联系的 PSG 也能产生并保护无能隙激发。由量子序产生并保护的无能隙激发可以是无能隙规范玻色子和/或无能隙费米子。

在本节中，我们将研究立方格子上一个具体的 $SU(N_f)$ 自旋模型（Wen, 2002a）。该模型的基态具有非平庸的量子序（即弦网凝聚）。结果，模型具有演生的无质量 $U(1)$ 规范玻色子（人工光）以及无质量的带电费米子（人工电子与质子）。换言之，$SU(N_f)$ 自旋模型具有一个演生的 QED！基于这种 $SU(N_f)$ 自旋模型的世界将与我们自己的世界十分相似；两个世界都由通过库仑相互作用耦合的带电费米子（即电子与质子）构成。更复杂的立方格子自旋模型甚至可以导致含有光子、胶子、轻子和夸克的演生 QED 与 QCD（Wen, 2003b）。因此，如果有一天我们发现自然界中所有基本粒子都从一个推广的自旋模型（即局域玻色模型）中演生出来，那也不足为奇。

\subsection{立方格子上的 $SU(N_f)$ 自旋模型}

该模型由 $SU(N_f)$ 自旋构成，立方格子的每个格点上放有一个自旋。这里 $N_f$ 为偶数，格点用三个整数标记，即 $\pmb i=(i_x,i_y,i_z)$。每个格点上的态
\begin{equation*}
|a_1,a_2,\dots,a_{N_f/2}\rangle,\qquad a_1,\dots,a_{N_f/2}=1,2,\dots,N_f
\end{equation*}
构成 $SU(N_f)$ 的秩为 $N_f/2$ 的反对称张量表示。每个格点上的自旋算符 $S^{a\bar b}_{\pmb i}$（$a,\bar b=1,2,\dots,N_f$）构成 $SU(N_f)$ 的伴随表示：
\begin{gather*}
[S^{a\bar b}_{\pmb i},S^{c\bar d}_{\pmb i}]=\delta^{\bar b c}S^{a\bar d}_{\pmb i}-\delta^{a\bar d}S^{c\bar b}_{\pmb i}\\
S^{a\bar b}_{\pmb i}\delta_{a\bar b}=0\\
(S^{a\bar b}_{\pmb i})^{\dagger}=S^{\bar b a}_{\pmb i}\tag{10.7.1}
\end{gather*}
其中重复指标意味着求和。

令
\begin{equation*}
W_{\pmb i}=-N_f^{-4}S^{ab}_{\pmb i+\pmb x}S^{bc}_{\pmb i+\pmb x+\pmb y}S^{cd}_{\pmb i+\pmb y}S^{da}_{\pmb i}
\end{equation*}
我们把 $SU(N_f)$ 模型的哈密顿量取为（Wen, 2002a）
\begin{equation*}
H=g\sum_{\pmb i}[W_{\pmb i}+h.c]\tag{10.7.2}
\end{equation*}

\subsection{$SU(N_f)$ 模型的基态}

为理解基态，让我们引入 $SU(N_f)$ 自旋模型的费米子表示（Affleck and Marston, 1988）。我们先引入构成 $SU(N_f)$ 基础表示的 $N_f$ 个费米子算符 $\psi_a$。每个格点上的希尔伯特空间于是由具有 $N_f/2$ 个费米子的态构成。自旋算符由下式给出：
\begin{equation*}
S^{a\bar b}=\psi_a\psi^{\dagger}_{\bar b}-\frac{1}{2}\delta_{a\bar b}
\end{equation*}
如果我们引入
\begin{equation*}
\hat\chi_{\pmb i\pmb j}=N_f^{-1}\psi^{\dagger}_{a\pmb i}\psi_{a\pmb j}
\end{equation*}
则 $W_{\pmb i}$ 由下式给出：
\begin{equation*}
W_{\pmb i}=\hat\chi_{\pmb i,\pmb i+\pmb x}\hat\chi_{\pmb i+\pmb x,\pmb i+\pmb x+\pmb y}\hat\chi_{\pmb i+\pmb x+\pmb y,\pmb i+\pmb y}\hat\chi_{\pmb i+\pmb y,\pmb i}+O(1/N_f)
\end{equation*}
利用费米子表示，我们可以把哈密顿量 (10.7.2) 写成
\begin{equation*}
H=g\sum_{\pmb p}\Big[\prod_{\square}\hat\chi_{\pmb i\pmb j}+h.c.\Big]+O(N_f^{-1})\tag{10.7.3}
\end{equation*}
其中 $\sum_{\pmb p}$ 对所有方格 $\pmb p$ 求和，$\prod_{\square}$ 对方格 $\pmb p$ 周围的四条链求积；也就是说，$\prod_{\square}\hat\chi_{\pmb i\pmb j}\allowbreak=\hat\chi_{\pmb i_1\pmb i_2}\allowbreak\hat\chi_{\pmb i_2\pmb i_3}\allowbreak\hat\chi_{\pmb i_3\pmb i_4}\allowbreak\hat\chi_{\pmb i_4\pmb i_1}$，其中 $\pmb i_1,\pmb i_2,\pmb i_3,\pmb i_4$ 是方格 $\pmb p$ 周围的四个格点。

我们注意到，在大 $N_f$ 极限下，$\hat\chi_{\pmb i\pmb j}$ 彼此对易，即 $[\hat\chi_{\pmb i\pmb j},\hat\chi_{\pmb i'\pmb j'}]=O(1/N_f)$。另外，$\hat\chi_{\pmb i\pmb j}$ 是 $N_f$ 项之和。尽管每一项都有很强的量子涨落，其和却几乎没有涨落。因此，在大 $N_f$ 极限下 $\hat\chi_{\pmb i\pmb j}$ 的行为像一个 c 数。也就是说，对一个给定的态，我们可以把 $\hat\chi_{\pmb i\pmb j}$ 换成一个复数 $\chi_{\pmb i\pmb j}$。这样的态的能量为 $E=g\sum_{\pmb p}[\prod_{\square}\chi_{\pmb i\pmb j}+h.c.]+O(N_f^{-1})$。不同的态有不同的 $\chi_{\pmb i\pmb j}$，因而有不同的能量。

$SU(N_f)$ 模型的基态可以这样得到：选取一组使 $E=g\sum_{\pmb p}[\prod_{\square}\chi_{\pmb i\pmb j}+h.c.]$ 极小的 $\chi_{\pmb i\pmb j}$。这里我们遇到一个困难：$\chi_{\pmb i\pmb j}$ 并不独立，我们不能任意选取它们的值。例如，我们不能让所有 $\chi_{\pmb i\pmb j}$ 同时为零。我们将用 9.1.1 小节中的投影构造\footnote{9.1.1 小节中的讨论针对的是 $N_f=2$ 的 $SU(N_f)$ 模型。}来克服这一困难。我们从平均场哈密顿量
\begin{equation*}
H_{\mathrm{mean}}=-\sum_{\langle\pmb i\pmb j\rangle}\tilde\chi^{\dagger}_{\pmb i\pmb j}\psi^{\dagger}_{\pmb i}\psi_{\pmb j}+\sum_{\pmb i}a_0(\pmb i)\psi^{\dagger}_{\pmb i}\psi_{\pmb i}
\end{equation*}
出发。我们这样选取 $a_0(\pmb i)$，使得 $H_{\mathrm{mean}}$ 的平均场基态，即 $|\Psi^{\{\tilde\chi_{\pmb i\pmb j}\}}_{\mathrm{mean}}\rangle$，平均每格点拥有 $N_f/2$ 个费米子。$SU(N_f)$ 模型的态可以由平均场态 $|\Psi^{\{\tilde\chi_{\pmb i\pmb j}\}}_{\mathrm{mean}}\rangle$ 投影到每格点恰有 $N_f/2$ 个费米子的子空间中来构造：
\begin{equation*}
|\Psi^{\{\tilde\chi_{\pmb i\pmb j}\}}\rangle=\mathcal P|\Psi^{\{\tilde\chi_{\pmb i\pmb j}\}}_{\mathrm{mean}}\rangle
\end{equation*}
现在我们可以选取 $\tilde\chi_{\pmb i\pmb j}$，使平均能量 $\langle\Psi^{\{\tilde\chi_{\pmb i\pmb j}\}}|H|\Psi^{\{\tilde\chi_{\pmb i\pmb j}\}}\rangle$ 极小，从而得到基态的尝试能量与尝试波函数。

在大 $N_f$ 极限下，投影 $\mathcal P$ 只使波函数 $|\Psi^{\{\tilde\chi_{\pmb i\pmb j}\}}_{\mathrm{mean}}\rangle$ 发生很小的改变。因此，为得到基态能量，我们可以转而做一件更容易的计算：求 $\langle\Psi^{\{\tilde\chi_{\pmb i\pmb j}\}}_{\mathrm{mean}}|H|\Psi^{\{\tilde\chi_{\pmb i\pmb j}\}}_{\mathrm{mean}}\rangle$ 的极小。

极小化按如下步骤进行。我们先挑出一组 $\tilde\chi_{\pmb i\pmb j}$，然后选取 $a_0(\pmb i)$ 使得 $\langle\Psi^{\{\tilde\chi_{\pmb i\pmb j}\}}_{\mathrm{mean}}|\psi^{\dagger}_{\pmb i}\psi_{\pmb i}|\Psi^{\{\tilde\chi_{\pmb i\pmb j}\}}_{\mathrm{mean}}\rangle=N_f/2$，接着计算 $\chi_{\pmb i\pmb j}=N_f^{-1}\langle\Psi^{\{\tilde\chi_{\pmb i\pmb j}\}}_{\mathrm{mean}}|\psi^{\dagger}_{\pmb i}\psi_{\pmb j}|\Psi^{\{\tilde\chi_{\pmb i\pmb j}\}}_{\mathrm{mean}}\rangle$。这些 $\chi_{\pmb i\pmb j}$ 就是上面讨论过的算符 $\hat\chi_{\pmb i\pmb j}$ 所对应的 c 数。我们看到，$\tilde\chi_{\pmb i\pmb j}$ 是独立的，其值可以任意选取。非独立的 $\chi_{\pmb i\pmb j}$ 是 $\tilde\chi_{\pmb i\pmb j}$ 的函数，即 $\chi_{\pmb i\pmb j}(\{\tilde\chi_{\pmb i\pmb j}\})$。基态通过在改变 $\tilde\chi_{\pmb i\pmb j}$ 的同时使 $E=g\sum_{\pmb p}[\prod_{\square}\chi_{\pmb i\pmb j}(\{\tilde\chi_{\pmb i\pmb j}\})+h.c.]$ 极小而得到。

我们发现，基态可以通过如下选取 $\tilde\chi_{\pmb i\pmb j}$ 得到：
\begin{align*}
\tilde\chi_{\pmb i,\pmb i+\hat{\pmb x}}&\equiv-\mathrm{i}, & \tilde\chi_{\pmb i,\pmb i+\hat{\pmb y}}&\equiv-\mathrm{i}(-)^{i_x},\\
\tilde\chi_{\pmb i,\pmb i+\hat{\pmb z}}&\equiv-\mathrm{i}(-)^{i_x+i_y}, & a_0(\pmb i)&=0
\end{align*}
它引导出如下一组 $\chi_{\pmb i\pmb j}$：
\begin{align*}
\bar\chi_{\pmb i,\pmb i+\hat{\pmb x}}&\equiv-\mathrm{i}|\chi|, & \bar\chi_{\pmb i,\pmb i+\hat{\pmb y}}&\equiv-\mathrm{i}(-)^{i_x}|\chi|,\\
\bar\chi_{\pmb i,\pmb i+\hat{\pmb z}}&\equiv-\mathrm{i}(-)^{i_x+i_y}|\chi| & a_0(\pmb i)&=0.\tag{10.7.4}
\end{align*}
拟设 $\bar\chi_{\pmb i\pmb j}$ 在每个方格中都有 $\pi$ 通量。结果是 $\prod_{\square}\chi_{\pmb i\pmb j}=-|\chi|^{4}$。正是这个负号使上述 $\chi_{\pmb i\pmb j}$ 极小化 $g\sum_{\pmb p}[\prod_{\square}\chi_{\pmb i\pmb j}+h.c.]$（假定 $g>0$）。

由拟设 $\bar\chi_{\pmb i\pmb j}$ 描写的基态不破坏任何对称性。然而，该基态具有非平庸的量子序。下面我们将证明，基态支持无质量的 $U(1)$ 规范玻色子与无质量的带电费米子。

\subsection{$SU(N_f)$ 模型的低能动力学}

$\chi_{\pmb i\pmb j}$ 的涨落描写我们 $SU(N_f)$ 自旋模型中的一个集体模式。这一集体模式的动力学由立方格子上的一个经典场论描写。$|\chi|$ 的涨落对应于一种有质量的激发，可以忽略。集体模式的低能涨落由围绕 $\bar\chi_{\pmb i\pmb j}$ 的相位涨落描写，即 $\chi_{\pmb i\pmb j}=\bar\chi_{\pmb i\pmb j}\mathrm{e}^{\mathrm{i}a_{\pmb i\pmb j}}$。如 9.1.1 小节所讨论的，由 $a_{\pmb i\pmb j}$ 描写的涨落是 $U(1)$ 规范涨落。$U(1)$ 规范涨落的物理波函数由变形拟设的投影平均场态给出：
\begin{equation*}
\mathcal P|\Psi^{\{\bar\chi_{\pmb i\pmb j}\mathrm{e}^{\mathrm{i}a_{\pmb i\pmb j}}\}}_{\mathrm{mean}}\rangle.\tag{10.7.5}
\end{equation*}

把上式代入式 (10.7.3)，就得到 $a_{\pmb i\pmb j}$ 的如下有效哈密顿量：
\begin{equation*}
H=-2g\sum_{\pmb p}|\chi|^{4}\cos(\Phi_{\pmb p})\tag{10.7.6}
\end{equation*}
其中
\begin{equation*}
\Phi_{\pmb p}=a_{\pmb i_1\pmb i_2}+a_{\pmb i_2\pmb i_3}+a_{\pmb i_3\pmb i_4}+a_{\pmb i_4\pmb i_1}
\end{equation*}
而 $(\pmb i_1,\pmb i_2,\pmb i_3,\pmb i_4)$ 是方格 $\pmb p$ 周围的四个格点。哈密顿量 (10.7.6) 描写一个 $U(1)$ 格点规范理论：$a_{\pmb i\pmb j}$ 是格点 $U(1)$ 规范场，$\Phi_{\pmb p}$ 是通过方格 $\pmb p$ 的 $U(1)$ 通量。

$U(1)$ 规范涨落 (10.7.5) 不是唯一类型的低能激发。第二类低能激发由带有若干粒子--空穴激发的投影平均场态给出，例如
\begin{equation*}
\mathcal P\psi^{a}_{\pmb i_1}\psi^{b\dagger}_{\pmb i_2}|\Psi^{\{\bar\chi_{\pmb i\pmb j}\}}_{\mathrm{mean}}\rangle.\tag{10.7.7}
\end{equation*}
这些激发对应于带电费米子。

为确定费米子激发的动力学，我们写出
\begin{equation*}
\hat\chi_{\pmb i\pmb j}=\bar\chi_{\pmb i\pmb j}+N_f^{-1}:\psi^{\dagger}_{a,\pmb i}\psi_{a,\pmb j}:\tag{10.7.8}
\end{equation*}
其中 $\bar\chi_{\pmb i\pmb j}=\langle\hat\chi_{\pmb i\pmb j}\rangle$。把式 (10.7.8) 代入式 (10.7.3)，我们得到
\begin{equation*}
H_f=gN_f^{-1}\sum_{\pmb i\pmb j}[\chi'_{\pmb i\pmb j}:\psi^{\dagger}_{\pmb j}\psi_{\pmb i}:+h.c.]+\sum_{\pmb i}a_0(\pmb i):\psi^{\dagger}_{\pmb i}\psi_{\pmb i}:+O\big((:\psi^{\dagger}\psi:)^{2}\big)\tag{10.7.9}
\end{equation*}
其中
\begin{gather*}
\chi'_{\pmb i,\pmb i+\hat{\pmb x}}\equiv4\mathrm{i}|\chi|^{3},\\
\chi'_{\pmb i,\pmb i+\hat{\pmb y}}\equiv4\mathrm{i}(-)^{i_x}|\chi|^{3},\\
\chi'_{\pmb i,\pmb i+\hat{\pmb z}}\equiv4\mathrm{i}(-)^{i_x+i_y}|\chi|^{3}.\tag{10.7.10}
\end{gather*}
我们看到，式 (10.7.9) 描写 $\pi$ 通量相中费米子的跳跃。把式 (10.7.6) 与 (10.7.9) 相结合，就得到费米子与规范激发的下述低能有效哈密顿量：
\begin{align*}
H_{eff}=gN_f^{-1}\sum_{\pmb i\pmb j}[\chi'_{\pmb i\pmb j}\mathrm{e}^{\mathrm{i}a_{\pmb i\pmb j}}:\psi^{\dagger}_{\pmb j}\psi_{\pmb i}:+h.c.]&+\sum_{\pmb i}a_0(\pmb i):\psi^{\dagger}_{\pmb i}\psi_{\pmb i}:\\
&-2g\sum_{\pmb p}|\chi|^{4}\cos(\Phi_{\pmb p})\tag{10.7.11}
\end{align*}
在连续极限下，$-2g\sum_{\pmb p}|\chi|^{4}\cos(\Phi_{\pmb p})$ 变为 $\int\mathrm{d}^{3}\pmb x\,(3l_0 g)|\chi|^{4}\pmb B^{2}$，其中 $l_0$ 是立方格子的晶格常数，而 $\pmb B$ 是
% ----（书页 491 末句延续到书页 492 顶部 "strength of the U(1) gauge field. Comparing this with..."。
% ---- 半句译文止于"而 $\pmb B$ 是"，请 ch10_e 以 %JOIN% 起头续译"该 $U(1)$ 规范场的磁场强度。
% ---- 把它与……"，勿重复勿遗漏。
